\documentclass[12pt]{article}

\usepackage[margin=1in]{geometry}

\usepackage{setspace}

\setstretch{1.2}

\begin{document}



\title{Proposal for a National Quantum Computing and Sensing Center}

\author{[Institution Consortium]}

\date{\today}

\maketitle



\section{Overarching Motivation}



Quantum technologies are emerging as a critical frontier of innovation with profound economic and security implications. Around the world, major nations are investing heavily in quantum computing, communication, and sensing, recognizing these as strategic technologies alongside artificial intelligence and 5G . The global race for quantum leadership is in full swing, with stakes high for technological and economic dominance . For example, the United States launched a National Quantum Initiative (NQI) in 2018, investing approximately $1.2billion over five years, yet this is dwarfed by China’s estimated $10–15billion commitment in the same period . Likewise, the European Union’s Quantum Flagship is a $1+ billion program to boost quantum R&D across member states . These investments underscore the urgency for national leadership in quantum science: quantum computing promises to solve problems intractable for classical supercomputers, quantum communications could enable fundamentally secure networks, and quantum sensors offer unparalleled precision for navigation and detection. Quantum technologies thus have the potential to transform industries, create new markets, and redefine national security capabilities .



Despite growing efforts, there remain strategic gaps that this proposed Center will address. First, a coordinated national effort is needed to bridge the gap between fundamental research and applied technology development. Currently, state-of-the-art quantum devices are too limited in scale and functionality to realize their full promise . No single institution can easily marshal the breadth of expertise or infrastructure required to overcome these challenges. Our Center will integrate a highly interdisciplinary team—spanning physics, computer science, engineering, materials science, and other fields—to pursue convergent research solutions. As noted in prior national initiatives, quantum advances require a “highly integrated interdisciplinary approach” because many hard challenges span multiple domains . By uniting researchers across disciplines under one umbrella, the Center will foster the co-design of quantum hardware, software, and applications in tandem, accelerating progress beyond what isolated efforts could achieve.



Secondly, the Center will fill critical gaps in infrastructure and workforce development. Governments have emphasized the need for “shared-use quantum testing facilities” and a “whole-of-nation” approach to quantum R&D . We lack domestic facilities at scale for quantum device fabrication, cryogenic testing, and quantum–classical computing integration. This Center’s state-of-the-art facilities (detailed in Section 3) will serve as national user infrastructure, accessible to academia and industry to spur innovation. Moreover, a shortage of skilled quantum scientists and engineers threatens to slow the field’s growth . Our Center will tackle this by training the next-generation quantum workforce through hands-on research and extensive education/outreach programs (see Section 5). In summary, establishing a National Quantum Computing and Sensing Center will ensure our country keeps pace with global competitors, translates scientific breakthroughs into practical technologies, and secures leadership in a domain poised to revolutionize economic prosperity and national security.



\section{Research Topics}



The Center’s research will be organized around several core themes that collectively address the most crucial challenges and opportunities in quantum computing and quantum sensing. Each theme is described below, with example topics of inquiry.



\subsection{Quantum Computing Architectures}



We will pursue multiple quantum computing hardware paradigms, leveraging the strengths of each and developing hybrid approaches. Leading qubit architectures under investigation include superconducting circuits, trapped-ion systems, photonic qubits, and emerging platforms (e.g. semiconductor quantum dots and neutral atoms). Superconducting qubits (used by organizations like IBM and Google) offer fast gate speeds and are amenable to microfabrication, while trapped ions provide long coherence times and high-fidelity operations, and photonic qubit systems promise natural connectivity and room-temperature operation. Rather than commit to a single technology, the Center will explore a modular and heterogeneous approach: connecting different types of quantum processors into one network to form a distributed quantum computer . This paradigm, wherein modest-size processors (e.g. tens of qubits each) are linked via entanglement links, can leverage each platform’s advantages while mitigating individual limitations . Research thrusts will include: improving qubit coherence and gate fidelities in each platform; developing microwave and optical interconnects to entangle superconducting qubits with optical photons or trapped ions; and designing a scalable modular architecture. By advancing multiple architectures in parallel, with a common interfacing framework, the Center aims to produce a quantum computing system that is greater than the sum of its parts – a network of quantum “nodes” capable of tackling problems beyond the reach of any single device.



\subsection{Quantum Error Correction and Algorithms}



Enabling useful, large-scale quantum computation will require surmounting the challenges of noise and decoherence that plague today’s quantum hardware. A major research thrust is quantum error correction (QEC): encoding logical qubits across many physical qubits to detect and correct errors in real time. We will investigate QEC codes (such as surface codes and bosonic codes) and fault-tolerant logic protocols tailored to our target hardware platforms. A key goal is to achieve and surpass the break-even point where a logical qubit’s error rate becomes lower than those of the underlying physical qubits. This involves both improving hardware stability and developing efficient decoding algorithms. Concurrently, the Center will research new quantum algorithms and applications, especially those optimized for near-term intermediate-scale quantum (NISQ) devices as well as algorithms designed for the error-corrected, at-scale quantum computers of the future. This includes quantum algorithms for cryptography (e.g. Shor’s algorithm for factoring, Grover’s search), for optimization and machine learning, and for quantum simulation of physical systems. We will place special emphasis on algorithms for quantum chemistry, materials science, and complex systems, which can demonstrate quantum advantage even with noisy hardware. The co-development of algorithms and hardware is critical; insights from algorithm requirements (e.g. qubit count, circuit depth, tolerance to error) will inform hardware design. Likewise, advances in QEC will feed back into algorithm design by expanding the class of computations that can be executed reliably. By tackling error correction head-on and innovating in algorithm design, the Center addresses the presently “unknown pathway” to a quantum computer that can outperform classical supercomputers in useful tasks.



\subsection{Quantum Networking and Communication}



Networking quantum devices via entanglement will unlock powerful new capabilities in both computing and communications. In this theme, we will research quantum networks that connect quantum processors and quantum sensors over local and long distances. One thrust is developing the hardware and protocols for a “quantum internet”: entanglement distribution across fiber optic links and free-space (satellite) links, enabling secure communication and distributed quantum information processing. Prior demonstrations have shown the feasibility of long-distance quantum key distribution (QKD) – for example, a 2,000 km fiber QKD network between Beijing and Shanghai, complemented by satellite links up to 2,600 km . Building on such advances, our Center will engineer metropolitan-scale quantum network testbeds. This involves single-photon sources and detectors, quantum memory devices for storing entangled states, and quantum repeaters that extend range by overcoming loss and decoherence in communication channels. We will also explore novel communication protocols that harness uniquely quantum phenomena, such as device-independent QKD (whose security is guaranteed even if devices are untrusted) and entanglement-based secure multi-party communication. A parallel thrust is distributed quantum computing: using networks to link smaller quantum computers into a larger virtual machine. By transmitting quantum states between nodes, we can perform distributed algorithms – for instance, secure distributed search or multi-party consensus tasks that are impossible classically . Research will include developing a full networking stack for quantum information, from low-level physical interfaces (wavelength conversion, routing of entangled photons) up to high-level software for allocating networked qubits and managing entanglement generation, swapping, and purification. The long-term vision is an internet-like infrastructure of quantum devices, enabling provably secure communications and distributed computing capabilities that dramatically surpass today’s networks in certain tasks . Our Center will ensure the nation remains at the forefront of this emerging quantum networking revolution.



\subsection{Quantum Sensors and Metrology}



Quantum sensing is perhaps the most mature quantum technology to date, yet it continues to offer new frontiers for research and application . Quantum sensors exploit quantum effects (such as superposition and entanglement) to measure physical quantities with extraordinary precision, often beyond classical limits. The Center will develop a suite of quantum sensors and metrological standards with broad impact. Example sensor modalities include: atomic clocks for time/frequency standards, atom interferometers for gravity and acceleration sensing, superconducting quantum interference devices (SQUIDs) and NV-center diamond magnetometers for detecting minuscule magnetic fields, and quantum optical sensors for imaging and communications. These sensors can measure time, position, rotation, electromagnetic fields, and other quantities with unprecedented exactness . Our research will push the sensitivity, stability, and field-deployability of such devices. One focus is quantum inertial sensors (gyroscopes and accelerometers using ultracold atoms or trapped ions) that could enable navigation in the absence of GPS, by dead-reckoning motion with quantum-enhanced accuracy . Another focus is quantum electromagnetic sensors that approach the standard quantum limit and can detect phenomena like electromagnetic signals of hidden nuclear materials or the neural magnetic fields of brain activity. We will also investigate quantum metrology techniques to improve standards for voltage, current, and other units by using quantum references (e.g. single-electron devices, quantum Hall effect devices). A key challenge is transitioning exquisite laboratory demonstrations into real-world instruments—today, many quantum sensors require laboratory conditions (vibration isolation, magnetic shielding, cryogenics). The Center’s engineering efforts will miniaturize and ruggedize these sensors for practical use, in some cases in partnership with agencies like DARPA that seek deployable quantum sensor technology. Ultimately, advancing quantum sensing and metrology not only provides powerful new tools for science and industry, but also strengthens national security by enabling capabilities such as precise navigation and remote detection that were previously unattainable with classical sensors .



\subsection{Applications in Materials, Chemistry, Biology, Navigation, and Defense}



Driving the above technical areas is a focus on real-world applications that solve pressing problems across various domains. Quantum science is inherently multidisciplinary, and our Center will maintain strong collaborative threads with experts in materials science, chemistry, biology, and defense to guide and apply our research. A flagship topic is quantum simulation for materials and chemistry: using quantum computers to model complex quantum systems that are beyond classical simulation, such as high-temperature superconductors, advanced battery materials, or reaction mechanisms in catalysis and drug molecules. A large-scale quantum computer could, for example, accurately simulate the electronic structure of complex molecules and materials, aiding the design of new pharmaceuticals and materials for energy harvesting . Even in the near term, hybrid quantum-classical algorithms (like variational eigensolvers) are being explored for tasks such as predicting chemical reaction rates and material properties with higher accuracy than classical methods. The Center’s researchers in chemistry and materials will work closely with the quantum algorithm teams to apply early quantum processors to problems like catalyst design for sustainable chemistry and discovery of novel functional materials.



In biology and medicine, quantum technologies promise improved diagnostics and understanding of complex biomolecules. Quantum sensors such as NV-diamond magnetometers could enable high-resolution NMR imaging of proteins or new types of biosensors that detect minute molecular signals. Likewise, as quantum computing hardware scales, it offers a revolutionary way to tackle drug discovery challenges. Instead of brute-force screening via classical computing (which struggles with the combinatorial complexity of molecular interactions), quantum computing can efficiently explore the vast space of molecular configurations and binding interactions . For example, quantum algorithms can precisely model how a drug molecule binds to a target protein (taking into account quantum effects in bonding and solvent interactions) far faster than classical simulations . This could optimize lead identification in drug development, reducing the time and cost to bring new therapies to market. Our Center’s application thrust in the life sciences will partner with pharmaceutical researchers to test quantum approaches on real drug design problems, as recent collaborations have begun to demonstrate .



Another critical set of applications lies in navigation and precision timing. Quantum clocks and inertial sensors (gyroscopes/accelerometers) can provide positioning and timing independent of GPS. This has major implications for both civilian and defense use. Precision navigation enabled by quantum sensors would allow, for instance, autonomous vehicles or aircraft to maintain accurate positioning in GPS-denied environments (underwater, indoors, or during jamming) . The Center will collaborate with aerospace and defense partners to develop and field-test quantum positioning systems, which leverage atomic interferometry to achieve navigation-grade performance.



In the realm of secure communications and defense, our research addresses both the offensive and defensive sides of quantum technology. On one hand, quantum computers in the future could break current public-key encryption—so our work on quantum-resistant algorithms (post-quantum cryptography) will support national cybersecurity by validating new encryption standards that can withstand quantum attacks . On the other hand, quantum communication techniques (QKD and beyond) offer fundamentally secure channels that adversaries cannot eavesdrop on without detection. The Center will implement and test quantum communication links with government and industry partners to demonstrate ultra-secure networks for military and financial data. Quantum sensing also opens up novel capabilities in defense: for example, gravimetric sensors can detect subterranean structures or hidden tunnels by sensing tiny changes in gravity, and quantum magnetic sensors could detect stealth aircraft or submarines via their magnetic or gravitational signatures . By developing such technologies, the Center can contribute to capabilities like underground resource mapping and long-range detection of threats . In summary, a unifying principle of our application-driven research is to ensure that breakthroughs in quantum computing and sensing directly translate into solutions for societal needs—be it accelerating the discovery of life-saving drugs, enabling resilient navigation and communication infrastructure, or protecting national security with new detection and encryption technologies.



\section{Technical Capabilities and Infrastructure}



Figure 1: Cryogenic quantum computing infrastructure example. A Bluefors KIDE dilution refrigerator system (rendering shown) provides the ultra-low temperatures (~10 mK) required for superconducting quantum processors . The proposed Center will house multiple such cryogenic systems alongside advanced fabrication and networking facilities to support quantum computing and sensing experiments.



A core strength of this proposal is the comprehensive suite of technical facilities we will establish to support cutting-edge research in quantum computing and sensing. The Center will be headquartered in a purpose-designed facility hosting specialized laboratories, fabrication cleanrooms, and collaborative workspaces. Key capabilities include:



Cryogenic Quantum Computing Laboratories: Many quantum computing platforms (such as superconducting qubits and some semiconductor qubits) require temperatures near absolute zero to operate, necessitating sophisticated cryogenic infrastructure . The Center will acquire multiple dilution refrigerators and cryostats capable of reaching millikelvin temperatures, each instrumented for qubit device testing with high-density microwave wiring, cryogenic amplifiers, and optical fiber feedthroughs as needed. These “cold labs” will be shielded against vibration and electromagnetic interference, providing the stable environment qubits need. To complement the low-temperature setups, we will also maintain laser/optics labs for trapped-ion and photonic quantum computing experiments. These labs will be equipped with high-stability laser systems (with frequency locking to atomic transitions), ultra-high vacuum chambers for ion traps or cold atom experiments, and single-photon detectors. By supporting multiple hardware modalities under one roof, the Center enables cross-platform comparisons and hybrid systems development (for instance, interfacing trapped-ion qubits with superconducting qubits via optical links).
Nanofabrication Cleanroom: The Center will include a state-of-the-art cleanroom facility for fabricating quantum devices and materials. This cleanroom (Class 100 or better in critical areas) will house electron-beam lithography, photolithography, thin-film deposition systems (for superconducting films, semiconductor heterostructures, etc.), plasma etching, and precision metrology tools. Such facilities are essential for producing superconducting qubit chips (with Josephson junctions and resonators), nanophotonic circuits for quantum light manipulation, solid-state defect centers (like NV centers in diamond), and custom sensor components. The cleanroom will operate as a national facility, allowing Center researchers and external collaborators to prototype new device designs rapidly in-house. In addition, materials characterization instruments (scanning electron microscope, atomic force microscope, etc.) will be available to analyze qubit materials and interfaces. By tightly integrating fabrication with testing, our researchers can iterate quickly on improvements to qubit coherence times, photonic chip efficiency, and sensor sensitivity.
High-Performance Computing and Quantum Simulation Cluster: To support both quantum computing research and classical simulation needs, the Center will deploy a high-performance computing cluster with specialized software for quantum simulation, control, and data analysis. This cluster will enable researchers to simulate quantum circuits and error-correction protocols (often requiring heavy classical computation) and to run classical optimization routines that partner with quantum algorithms (for example, variational algorithm feedback loops). Furthermore, we plan to integrate a quantum computer testbed with the HPC environment to create a quantum-classical hybrid computing testbed. In collaboration with computing centers, one of the Center’s quantum processors (once operational) will be co-located or networked with a supercomputer to enable fast classical/quantum co-processing. This follows the successful model demonstrated in Finland, where a quantum computer (Helmi) was connected to the LUMI supercomputer to provide a hybrid HPC-Quantum service . Our hybrid testbed will allow researchers nationwide to run experiments that harness the strength of classical supercomputing for data handling and the power of quantum co-processors for the parts of problems where quantum speedup is possible.
Quantum Networking Testbed: The Center will establish an experimental quantum network to link quantum devices within the facility (and potentially to external nodes, such as a nearby university or national lab). This will include tens of kilometers of spooled low-loss optical fiber for testing fiber-based entanglement distribution, as well as free-space optical terminals on the facility’s rooftop for line-of-sight quantum communication (useful for experimenting with satellite quantum link technologies). We will deploy sources of entangled photon pairs, single-photon detectors, and timing synchronization systems to create a functional quantum link. This testbed will allow researchers to develop and prototype quantum network hardware (quantum repeaters, memory devices) and to run quantum communication protocols in a realistic environment. The eventual goal is to serve as a regional hub that can connect to larger quantum network initiatives nationally.
Quantum Sensor and Metrology Labs: In addition to computing, the Center will dedicate laboratory space to developing quantum sensors. For atomic/ionic sensors (like atom interferometers for gravity or atom-based magnetometers), we will have labs with magnetic shielding chambers, vacuum systems, laser cooling setups, and precision rotation/translation stages to simulate various conditions. For solid-state sensors (NV centers, spin defects, etc.), we will have low-temperature optical microscopy setups and microwave electronics for controlling spin qubits in materials. An anechoic chamber and EMI shielded room will support testing of sensors in low-noise electromagnetic conditions. The Center will also maintain a timing laboratory with active hydrogen masers and optical clocks, serving as a testbed for next-generation atomic clock comparisons and perhaps contributing to national time standards.




Collectively, these technical capabilities will make the Center a one-stop shop for quantum R&D, from device fabrication to complete system integration. Equally important, the infrastructure will be open to collaborators and industry partners. By providing shared-use facilities , we lower barriers to entry for innovative ideas and ensure that the broader quantum community benefits from the world-class equipment. In summary, the proposed Center’s infrastructure investment will establish a national asset: a collaborative environment with cutting-edge experimental capabilities that will accelerate quantum technology development and demonstration.



\section{Applications and Impact}



A driving rationale for this Center is to deliver tangible benefits by applying quantum technologies to solve important real-world problems. We have identified several high-impact application domains where breakthroughs are anticipated:



\paragraph{Computational Chemistry and Drug Discovery:} Quantum computing offers a new paradigm for simulating molecular and materials systems with high accuracy, which could vastly improve the process of discovering new drugs and chemicals. In classical drug discovery, predicting how a drug molecule will bind to a target protein and how molecules interact involves extreme computational complexity, often requiring approximate models. Quantum computers can in principle model these quantum-mechanical interactions exactly for larger, more complex molecules than classical computers can handle. By leveraging quantum algorithms, researchers can explore vast chemical reaction spaces more efficiently . For example, simulating enzyme mechanisms or drug-receptor binding affinities with quantum algorithms could identify promising drug candidates much faster than current methods . This has direct societal impact: faster development of medications and novel materials (for energy storage, carbon capture, etc.). The Center’s collaborative team including chemists and quantum algorithm experts will target demonstrator projects such as calculating the properties of complex biomolecular systems (protein-ligand interactions, protein folding pathways) on our quantum platforms, in partnership with pharmaceutical companies. Early successes have already been noted in industry collaborations where quantum computers were used to analyze protein hydration or binding sites , and we aim to push these frontiers further. Achieving quantum-enabled drug discovery would strengthen healthcare by reducing R&D costs and enabling more precise, effective therapies.



\paragraph{Secure Communications and Data Security:} Ensuring secure information exchange is a cornerstone of national security and digital economy. Quantum technology impacts secure communications in two ways: enabling new secure communication methods and threatening classical encryption. On the defensive side, \textbf{quantum key distribution} (QKD) allows two parties to share encryption keys with security guaranteed by quantum physics – any eavesdropping on the quantum channel can be detected. Our Center will advance QKD and more general entanglement-based networking as described in Section 2.3, with the goal of deploying pilot secure communication links (for example, between government research sites or between a bank and its data center). We will also explore integrating QKD with conventional networks and developing quantum network simulators to test scalability. On the offensive side, a full-scale quantum computer could break widely used public-key cryptosystems (RSA, ECC) via Shor’s algorithm. Anticipating this, the world is transitioning to \textbf{post-quantum cryptography} (PQC) – classical encryption algorithms that are resistant to quantum attacks. Our Center’s researchers in quantum algorithms and mathematics will contribute to the analysis and implementation of PQC, helping to ensure that new standards (like those selected by NIST) are robust and efficiently deployable . In tandem, by providing quantum testbeds, we allow agencies and companies to empirically test the security of both quantum and classical cryptographic schemes in realistic conditions. The outcome will be communication infrastructure that can withstand future quantum threats and harness quantum technology for unconditional security where needed.



\paragraph{Precision Navigation and Timing (PNT):} Modern society and military operations rely heavily on the Global Positioning System (GPS) for navigation and timing. However, GPS signals are weak and vulnerable to jamming or spoofing. Quantum-based PNT systems can provide an alternative or complement to GPS, offering accurate navigation and timing that is immune to external interference. For instance, an aircraft or submarine equipped with a quantum inertial navigation system (using atom interferometers or cold-atom gyroscopes) could maintain precise course and timing for long periods without GPS input . The Center will collaborate with aerospace partners to develop such quantum inertial navigation units, aiming to demonstrate field-worthy prototypes. Similarly, portable atomic clocks with stability far beyond GPS clocks could be used to timestamp financial transactions or synchronize telecom networks with unprecedented precision. Our work in developing robust, compact atomic clocks (e.g. using optical clock transitions in ions or atoms) will support the creation of an improved timing backbone for critical infrastructure. The broader impact of quantum PNT technology includes safer autonomous vehicles (with quantum IMUs ensuring reliability if GPS fails), improved geophysical exploration (through gravity gradiometers mapping the underground), and enhanced resilience of military platforms to electronic warfare. By investing in quantum navigation and timing, the Center will help secure our transportation, communication, and defense systems against disruption.



\paragraph{Sensing and Detection Technologies:} Quantum sensors stand to revolutionize how we detect and monitor our environment, from the microscopic to the geophysical scale. One application under development is quantum-enhanced imaging: for example, using squeezed light or entangled photons in imaging systems to surpass classical resolution or sensitivity limits, which could improve biomedical imaging or remote sensing. Another critical application is the detection of what was previously undetectable. In defense and security, quantum magnetometers and gravimeters could be game-changers. They might detect stealth aircraft by the minute disturbances in the Earth’s magnetic field or gravitational field caused by the aircraft, or detect buried hazardous materials and nuclear substances by their subtle gravitational anomalies . The Center’s quantum sensing research (Section 2.4) directly feeds these applications. We will work with defense labs to test quantum sensors’ ability to detect underground tunnels, submarines, or hidden objects at stand-off distances. In the scientific realm, quantum sensors will enable new discoveries: for instance, arrays of atomic magnetometers could map brain activity with high spatial resolution, advancing neuroscience, or they could detect signals of elusive phenomena like dark matter . By developing more sensitive and robust quantum sensors, the Center will broaden humanity’s ability to observe and interact with the world, with outcomes ranging from improved medical diagnostics and environmental monitoring to enhanced national security through better threat detection .



Each of the above application areas benefits not only from the eventual quantum technologies themselves but also from cross-pollination with other sectors. Through our Education and Outreach (Section 5) and Partnership (Section 6) activities, we will ensure that end-users—whether they are biologists, clinicians, navigators, or defense analysts—are engaged in the research process. By doing so, the Center will remain focused on delivering solutions that are relevant and readily transitioned to practical use. The measure of success for the Center will ultimately be in these real-world impacts: new drugs discovered, communication networks secured, navigation systems hardened, and sensors deployed that together enhance technological leadership and societal well-being.



\section{Education and Outreach}



Training a quantum-capable workforce and engaging the public are central objectives of the proposed Center. Quantum science and engineering are fields that require highly specialized skills, as well as a broad understanding across physics, engineering, and computer science. Moreover, to fully harness quantum technologies for societal benefit, we must cultivate talent from all backgrounds and inform the public about quantum’s potential. The Center will implement a comprehensive Education and Outreach program with several components:



Workforce Development and Academic Programs: We will establish new academic programs and curricula to produce graduates fluent in quantum information science and technology (QIST). This includes creating an interdisciplinary \emph{Quantum Science and Engineering} graduate program (M.S. and Ph.D. levels) that draws on faculty from physics, electrical engineering, computer science, and related departments. Students will receive a core curriculum covering quantum mechanics, quantum computing algorithms, quantum hardware, and quantum communication, supplemented by electives in their specialization (e.g. photonics, theoretical computer science, etc.). Importantly, students will gain hands-on experience through research rotations in the Center’s labs, working on real experiments and devices. At the undergraduate level, we will develop a QIST minor or certificate program open to students in science and engineering, ensuring exposure to quantum topics early in the educational pipeline. The Center will also host an annual \textbf{Quantum Summer School} for advanced undergraduates and beginning graduate students nationwide, to accelerate their introduction to the field.



K-12 and Early Education Outreach: Recognizing that inspiration for STEM careers often begins early, the Center will actively engage with K-12 schools. We will partner with local school districts to incorporate basic quantum science demos and modules into their curricula. For example, we will develop age-appropriate teaching materials that introduce high school students to the concepts of qubits and cryptography, or demonstrate quantum sensor principles (like showing magnetic resonance imaging concepts using classroom experiments). The Center plans to host on-site tours and interactive lab days for high school classes, giving students a first-hand look at quantum experiments. It is crucial to include quantum concepts at early stages of education; quantum should be introduced at the K-12 level so that students build intuition before college . By creating an accessible and captivating set of educational resources (including perhaps a “Quantum in a Box” kit that can be loaned to schools), we will help sow the seeds for the next generation of quantum scientists.



Broadening Participation and Diversity: A vibrant and innovative quantum workforce must draw talent from all sectors of society. The Center is committed to diversity and inclusion in all its programs. We will proactively reach out to underrepresented groups in STEM (including women, certain minorities, and economically disadvantaged communities) to encourage their participation in quantum science. This involves forging partnerships with minority-serving institutions and community colleges to create internship pipelines into our labs. We will offer summer research internships for students from these partner institutions, providing mentorship and stipends to reduce barriers. Furthermore, the Center will support quantum workshops and bootcamps for professionals and students from backgrounds underrepresented in tech, to build confidence and skills. As highlighted in national policy discussions, reaching “underemployed and underserved” populations with quantum education is essential to tap an untapped talent pool . Our outreach will include programs in vocational and technical colleges to teach technicians the specialized skills (like cryogenic engineering, precision instrumentation) needed in quantum R&D, thereby expanding the workforce beyond just Ph.D. scientists. We will track our diversity and inclusion progress with metrics (e.g. internship demographics, retention rates) and continually improve our strategies.



Public Engagement and Quantum Awareness: Beyond formal education, the Center will serve as a focal point for public communication about quantum science and its implications. We plan to host public lectures, Q&A sessions, and annual open-house events where community members can tour our facilities and see demonstrations (such as viewing a real quantum computer or quantum levitation demonstrations). Especially as quantum technologies can sound mysterious or abstract, it is important to demystify them and articulate their benefits and risks to society. The Center will work with science museums and media to create exhibits or documentaries highlighting quantum research. In coordination with national efforts, we may contribute to a public awareness campaign to improve “quantum literacy” so that the general population and policy-makers are informed about what quantum technology is and isn’t . This could include contributing expert speakers to local science festivals, maintaining an engaging website with explainer articles and videos, and leveraging social media for outreach (e.g. a YouTube series featuring our students explaining quantum experiments). Our goal is to make quantum science accessible and exciting, inspiring public support and future student interest.



Through these Education and Outreach initiatives, the Center will address the widely recognized workforce challenge in quantum technologies . By increasing the number of trained quantum Ph.D.s and engineers, equipping workers with quantum-related skills, and engaging young students, we ensure a robust talent pipeline for the nation’s quantum ecosystem. Moreover, by emphasizing inclusivity and public engagement, we aim to build a quantum community that is not only technically skilled but also diverse, inclusive, and well-understood by the society it serves.



\section{Partnerships and Collaborations}



No single institution can achieve the ambitious goals of this proposal alone. A key strength of the Center will be its extensive network of partnerships spanning academia, industry, government, and international collaborators. By fostering collaborations, we will leverage complementary expertise, share resources, and accelerate innovation in quantum computing and sensing. Our partnership strategy includes:



Academic Consortia: The Center is proposed as a multi-institutional effort uniting several leading universities and research institutes. Each academic partner brings unique strengths—be it theoretical quantum computing, cryogenic hardware expertise, photonics, or domain applications like chemistry and materials. For example, University A (with a top-ranked physics program) will lead research on superconducting qubit hardware and fundamental algorithms, while University B (known for photonics) will spearhead integrated photonic quantum circuit development, and University C (with a strong engineering school) will focus on quantum sensor prototyping and deployment. We will formalize these collaborations via sub-awards and a shared governance structure for the Center. Regular all-hands meetings, cross-institution student exchanges, and shared facilities access will knit the academic partners into one cohesive team. This leveraging of diverse talent and facilities is necessary to cover the breadth of our research topics. The National Science Foundation’s establishment of Quantum Leap Challenge Institutes and similar centers has shown the value of such academic consortia for quantum research . Our Center will extend that model, creating a hub-and-spoke network that amplifies each partner’s contributions.



National Laboratories and Government Agencies: We will closely collaborate with U.S. national laboratories (and/or analogous national research entities, depending on country context) that have significant programs in quantum science. For instance, we anticipate partnerships with Department of Energy (DOE) labs such as Argonne, Oak Ridge, and NIST (National Institute of Standards and Technology). These labs offer world-class facilities (nanofabrication centers, characterization tools like advanced microscopy and cryogenic test facilities, etc.) and expertise that complement the academic environment. A concrete example is partnering with NIST on quantum metrology standards—our Center’s precision measurement experiments in quantum clocks could be done in coordination with NIST’s time and frequency division. We will also seek to participate in DOE’s existing quantum research centers and user facilities, plugging our researchers into broader national projects. Government agencies such as DARPA, IARPA, and NASA are potential sponsors/collaborators for more applied challenges (DARPA for defense-oriented quantum sensors, NASA for quantum navigation in space, etc.). Early engagement with such agencies will ensure our research aligns with mission needs and facilitates tech transfer into operational uses. Overall, cross-government collaboration will amplify impact ; the Center can serve as a bridge connecting fundamental research to the specific requirements of national security, space, and standards missions.



Industry Partnerships: Active involvement of industry partners is crucial for ensuring that our research leads to practical technologies and for tapping into additional expertise and resources. The Center has secured commitments from key industry players in the quantum ecosystem, including established corporations and startups. These include quantum hardware companies (manufacturers of quantum computers or components), software startups developing quantum algorithms, tech giants investing in quantum cloud services, and companies in sectors like pharmaceuticals, finance, and aerospace interested in quantum applications. We will form an Industry Advisory Board where our partners can provide guidance on research directions, share real-world problem sets for us to tackle, and facilitate commercialization of our innovations. For example, a partnership with a quantum computing company (like IBM, Google or a startup) might grant our researchers early access to cutting-edge prototype processors, or allow joint development of control electronics. An aerospace company might co-develop a quantum sensor and host field tests on their platform. Many industry partners have also expressed willingness to host student interns, which will enhance workforce development by giving students exposure to industrial R&D. By working closely with industry, we can identify and overcome practical engineering challenges early, increasing the technology readiness level (TRL) of our research. Furthermore, as emphasized by leaders in the field, close collaboration with industry helps “bring technologies to practical application” and speed up the transition from lab to market . We will protect intellectual property through the standard mechanisms (joint IP agreements, etc.) but with a philosophy of openness that encourages rapid knowledge transfer and possibly open-source contributions in areas like quantum software.



International Collaboration: Quantum science is a global endeavor, and international partnerships will enhance our Center’s capabilities while fostering a cooperative global quantum ecosystem. We plan to collaborate with leading quantum research centers and consortia abroad (for example, the EU Quantum Flagship projects, Canada’s Quantum institutes, Australia’s CQC2T, and others). These collaborations might take the form of student and researcher exchanges, jointly organized workshops, and sharing of data or access to unique facilities. By engaging like-minded nations in joint projects, we can pool resources and knowledge, accelerating progress for all . For instance, if a partner in Europe has a specialized quantum network testbed or a unique material for qubits, and our Center has complementary technology, sharing and benchmarking across both can be mutually beneficial. We will also participate in international standardization efforts for quantum technologies, ensuring that our work contributes to globally interoperable systems (e.g. standards for quantum cryptography protocols or qubit interface specifications). Given the importance of quantum tech in the coming decades, maintaining strong alliances through science diplomacy is also strategically important . Our Center will be a platform to strengthen these ties—hosting international visiting scholars, contributing to international conferences, and aligning our research objectives with global best practices.



Finally, the Center will maintain an internal structure to coordinate all these partnerships effectively. We will have dedicated coordinators or liaisons for each sector (academic, lab, industry, international) who report to the Center’s leadership. Regular meetings of a Partnerships Council will review collaboration outcomes and identify new opportunities. By actively managing our network of collaborators, we ensure that this large-scale initiative remains integrated and goal-oriented.



In conclusion, the establishment of this National Quantum Computing and Sensing Center will catalyze a powerful consortium of talents and resources. Through strategic partnerships, we multiply our strengths: academic excellence, industrial know-how, government mission focus, and international perspectives all coalesce in this Center. This collaborative approach not only increases the likelihood of scientific breakthroughs (by bringing many minds to the task), but also ensures that when breakthroughs occur, they swiftly make their way into real-world use. In doing so, the Center will position our nation at the forefront of the quantum revolution, leading through unity and shared purpose in this critical scientific endeavor.



\vfill\noindent {\footnotesize\textbf{References and Notes:} Citations marked with 【†】 refer to source materials and strategic documents that informed this proposal. Full references are available upon request.}



\end{document}
